% Expanded structural and analytic verification material, 29 September 2026.
% This file is an additive fragment. It never changes the original manuscript.
% All calculations below are transcriptions or elementary deductions from
% archived results; no numerical research producer was run for this expansion.

\section{Additional structure: an atlas of the commuting algebras}
\label{expstructure:atlas}

Several different commutants occur in the construction. Their central
charges and their ambient algebras must be specified together. This
section collects the scattered structural calculations, including useful
routes that the main proof does not need. The exact inputs have the same
status as the finite certificate inputs in the main text. The additional
all-weight arguments below are reproduced from the named research notes
for verification; they are not new numerical replays or independent
external validation. In particular, identifying an ambient commutant
does not identify its intersection with the actual generated algebra
$U$.

We use $\mathcal T$ for the lattice--Potts core of charge $14/5$,
$P=K(\mathfrak{sl}_2,3)$, $Q=K(\mathfrak{sl}_2,13)$, and $R=A_L$
for the order-thirteen extension of $\mathcal T\otimes Q$. Write
$s_R=\omega_B-\omega_R$, so $c(s_R)=18/5$. This notation avoids
confusing the core $\mathcal T$ with the translation operator $D$.
The three ambient algebras remain
$U\subset B=V_{E_6\oplus A_2}\subset H=V_{E_8}$.

\begin{longtable}{@{}p{0.28\linewidth}p{0.12\linewidth}p{0.53\linewidth}@{}}
\toprule
Algebra or stress&Charge&Information available\\
\midrule\endhead
$P$&$4/5$&Actual Potts embedding, including its weight-three generator.\\
$\mathcal T$&$14/5$&Actual order-three extension of a rank-two lattice
and $P$; four currents.\\
$Q$&$8/5$&Actual parafermion embedding obtained from complete finite
polynomial OPE rows.\\
$R=A_L$&$22/5$&Actual order-thirteen simple-current extension;
fourteen simple modules and four currents.\\
$\Com(\mathcal T,B)$&$26/5$&An ambient rational frame
$V_M\otimes L(1/2,0)\otimes L(7/10,0)$.\\
$C_B=\Com(R,B)$&$18/5$&Five weight-two states; supplementary
all-weight identification with the $(4,8,1;13)$ diagonal coset.\\
$C_H=\Com(R,H)$&$18/5$&Thirteen weight-two states; supplementary
all-weight identification with an $F_4$ diagonal coset.\\
$C_U=\Com(R,U)$&$18/5$&Two weight-two states and two weight-three
states under the properness and finite-grade conclusions. No complete
character is supplied here.\\
$C_6=\Com(V_L,U)$&$6$&Lattice double commutant; supplementary
regularity and category factorization; seven weight-two states.\\
\bottomrule
\end{longtable}

The charge-$26/5$ and charge-$18/5$ ambient cosets are different
objects. The two charge-$18/5$ ambient cosets are also different from
$C_U$. Finally, the charge-$6$ lattice coset removes only the Cartan
plane; it retains the Potts and parafermion factors. These distinctions
explain several otherwise incompatible-looking dimensions in the old
records.

\subsection{The determinant-thirty-nine lattice and its actual extension}
\label{expstructure:lattice-core}

In the simple-root coordinates of the construction put
\[
\begin{split}
a_0&=(-1,-2,-3,-2,-1,-2,0,0),\\
H_R&=(5,8,9,8,5,0,2,2),\\
t_0&=(1,1,0,1,1,-3,1,1)=\tfrac12(3a_0+H_R).
\end{split}
\]
Their pairings are
$a_0^2=2$, $H_R^2=78$, $(a_0,H_R)=0$, $t_0^2=24$ and
$(a_0,t_0)=3$. Thus
\[
 L=\ZZ a_0+\ZZ t_0,
 \qquad G_L=\begin{pmatrix}2&3\\3&24\end{pmatrix},
 \qquad \det G_L=39.
\]
The vector $r_0=t_0/3=a_0/2+H_R/6$ has norm $8/3$. The exact
generating field $u$ has Potts weight $2/3$ and plane-Heisenberg
weight $4/3$. The identity
\begin{equation}
 u_{-5}(u_{-3}u)=-\frac35e^{t_0}
 \label{expstructure:triple-product}
\end{equation}
places the lattice field $e^{t_0}$ in $U$. PCT gives $e^{-t_0}$,
and the left affine factor already supplies $e^{\pm a_0}$ and
the required Heisenberg directions. This is the actual field-level
reason for $V_L\subset U$; it is not inferred from a character.
The inner product formula fixes the leading power in a product of
lattice exponentials. For the triple product in
\eqref{expstructure:triple-product}, only the leading terms at the
specified powers contribute. A basis of the full weight-twelve
ambient space is therefore unnecessary.

Let $\psi$ be the Potts simple current of weight $2/3$. The extension
is
\begin{equation}
 \mathcal T=(V_L\otimes P)
 \oplus(V_{L+r_0}\otimes\psi)
 \oplus(V_{L+2r_0}\otimes\psi^2).
 \label{expstructure:T-extension}
\end{equation}
The order-three support is closed under fusion. In the existing
unitary ambient algebra, nonzero multiplication with the conjugate
sector has a nonzero vacuum pairing, so multiplication by a supported
invertible sector is invertible. Connectedness then gives multiplicity
one. The commutant-energy argument removes a possible hidden
positive-energy complementary factor. These observations promote the
three witnessed charge sectors to the actual embedded extension.

The core stress and its complement are
\[
 \omega_{\mathcal T}
 =\frac14h_{a_0}(-1)^2\one
   +\frac1{156}h_{H_R}(-1)^2\one+p,
 \qquad c_{\mathcal T}=2+\frac45=\frac{14}5.
\]
Here the two Heisenberg terms have central charge one each. The
source is \artifact{hi_cft_20260923_proof_rational_core.json}, whose
exact record includes the three terms of $u_{-3}u$, the triple-product
coefficient, the lattice Gram matrix and the stress commutators.
Its historical caveat that this core alone does not prove rationality
of $U$ remains correct; the positive-stress proof uses additional
factors.

\subsection{Discriminant conventions and why two quadratic forms agree}

The discriminant generator $g=H_R/39$ has
$g^2/2=1/39$. Its pairings with the lattice basis are $0,1$, so it
lies in $L^*$, and its order is $39$. The class of $r_0$ is $26g$:
indeed $r_0-26g=a_0/2-H_R/2=2a_0-t_0\in L$.
Condensing the order-three diagonal current in
\eqref{expstructure:T-extension} leaves the order-thirteen pointed
generator $3g$, with quadratic form
\[
 q(j_{\rm old})=\frac{3j_{\rm old}^2}{13}\pmod1.
\]
The top-bound and core-extension files instead use the actual right
Cartan charge $j=6j_{\rm old}\pmod{13}$. Since $6^{-1}=11$ and
$3\cdot11^2\equiv-1\pmod{13}$, the same form becomes
\[
 q(j)=-\frac{j^2}{13}\pmod1.
\]
Thus the two formulas in the archive are a change of labels, not a
disagreement of discriminant forms. The remaining category is
\[
 \Rep(\mathcal T)\simeq
 \mathrm{Fib}\boxtimes\mathcal C(\ZZ/13,-j^2/13),
 \qquad
 \operatorname{Dim}\Rep(\mathcal T)
 =13\frac{5+\sqrt5}{2}.
\]
The equality of this core's central charge with that of $G_{2,1}$
does not identify the VOAs: $\mathcal T_1$ has dimension four, while
the affine $G_2$ vacuum algebra has fourteen currents.

\subsection{An explicit rational frame in the larger ambient commutant}
\label{expstructure:ambient26}

Let $\alpha=b-a$, $\beta=c-b$, and $\gamma=\alpha+\beta$, using
the three momenta in the original $u$. The Gram matrix is
$\left(\begin{smallmatrix}4&-2\\-2&4\end{smallmatrix}\right)$.
If $p_\delta$ denotes the exponential coefficient of the actual
Potts stress, set
\[
 e_I=\frac1{16}h_\alpha(-1)^2\one
       -\frac54\bigl(p_\alpha e^\alpha+p_{-\alpha}e^{-\alpha}\bigr),
 \qquad
 e_{TI}=\omega_{\sqrt2A_2}-p-e_I.
\]
The saved exact calculations give mutually commuting conformal
vectors with charges $1/2$ and $7/10$. In the unitary realization
they generate the simple Ising and tricritical Ising Virasoro VOAs.
The rank-four perpendicular lattice has basis
\[
\begin{aligned}
m_1&=(-1,1,0,0,0,0,0,0),&m_2&=(0,0,0,-1,1,0,0,0),\\
m_3&=(0,0,-1,0,0,0,1,0),&m_4&=(0,0,-1,0,0,0,0,1)
\end{aligned}
\]
and Gram matrix
\[
G_M=\begin{pmatrix}
6&0&1&1\\0&6&-1&-1\\1&-1&4&1\\1&-1&1&4
\end{pmatrix}.
\]
The lattice $L\oplus\sqrt2A_2\oplus M$ has index $156$ in
$E_6\oplus A_2$. Consequently the actual ambient commutant contains
\[
 C_0=V_M\otimes L(1/2,0)\otimes L(7/10,0),
 \qquad \omega_{C_0}=\omega_B-\omega_{\mathcal T}.
\]
The charge check is $4+1/2+7/10=26/5$. Because the conformal vector
is shared, restrictions to this regular frame have no independently
shifted copies of its simple modules: every copy starts at that
module's prescribed lowest weight. Finite-dimensional ambient graded
spaces then make the extension finite. The archived argument applies
the unitary extension theorem to obtain strong rationality of
$\Com(\mathcal T,B)$.

The exact record
\artifact{hi_cft_20260923_proof_ambient_coset.json} also explicitly
reports $e_I\notin S_2$ and $e_{TI}\notin S_2$. It therefore does
not give this frame inside $\Com(\mathcal T,U)$. The older numerical
record \artifact{hi_cft_20260923_proof_core_coset.json} found dimensions
$4$ and $15$ in the intersections with the selected $S_2,S_3$, with
a Potts-kernel residual approximately $7.665\times10^{-11}$.
Those two numerical dimensions are retained only as historical finite
span observations. Neither is used here as a certified full coset
character.

\subsection{Order-thirteen extension: the ten additional weight-two states}
\label{expstructure:R-extension}

The zero-$s_R$-energy part of the exact $120$-dimensional $S_2$ has
dimension $25$, while $(\mathcal T\otimes Q)_2$ has dimension $15$.
The ten additional states separate into the following simultaneous
energy types:
\[
\begin{array}{c|c|c|c|c}
\text{number}&(H_L,H_R)&h_{\mathcal T}&h_Q&h_P\\\hline
6&(\pm1,\pm5),(0,\pm8)&14/13&12/13&2/3\\
4&(\pm1,\pm3)&4/13&22/13&\text{not needed for the witness}
\end{array}
\]
In the first row, the four signs in $(\pm1,\pm5)$ are independent.
The saved vectors are actual joint eigenvectors in $S_2$. For charge
$j=5$, the Fibonacci-$\tau$ core module begins at $31/65$;
$14/13-31/65=3/5$ is not integral. Thus the witnessed core factor
has Fibonacci label $1$. The actual parafermion lowest-weight
formula singles out the two inverse simple currents at $12/13$.
By contrast, $22/13$ alone does not distinguish a simple current
from the $i=10$ parafermion sector. The second row is a consistency
check, not the step forcing the extension.

The pointed form of $\Rep(\mathcal T\otimes Q)$ is
$-(j^2+k^2)/13$. Its two maximal isotropic lines are $k=\pm5j$.
The witnessed nonzero pair generates one full line; both lines
cannot coexist in an integer-graded local extension since they
generate the entire pointed group, including nonintegral twists.
After reversing the parafermion charge convention if necessary,
\[
 R=\bigoplus_{j\in\ZZ/13}\mathcal T_{1,j}\otimes Q_{0,-5j}.
\]
There is no fixed-point resolution: the order-thirteen action on
the charge pairs is free. The remaining simple labels are
$x_a=(1,a)$ and $y_a=(\tau,a)$, $0\leq a\leq6$, with the second
entry corresponding to even affine label $2a$. The category is
$\mathrm{Fib}\boxtimes\mathrm{SO}(3)_{13}$.

The thirteen vacuum-extension sectors have lowest weights
\[
\begin{array}{c|rrrrrrr}
 |j|&0&1&2&3&4&5&6\\\hline
 h&0&4&3&2&4&2&3
\end{array}
\]
where positive $|j|$ denotes two conjugate sectors. Only the charge
$\pm3$ and $\pm5$ sectors add weight-two states. This supplies an
independent explanation of $\dim R_1=4$ and $\dim R_2=25$.
The source is
\artifact{hi_cft_20260923_next_structure_order13_core.md} and its
exact spectral record
\artifact{hi_cft_20260923_next_structure_core_extension.json}.

\subsection{Why the core is already its double commutant}
\label{expstructure:core-maximality}

The actual lowest-weight table has only two integer-weight labels:
$x_0$ of weight $0$ and $x_5$ of weight $1$. Hence a connected
integer-graded conformal extension $D$ of $R$, with finite graded
spaces, has underlying module
\[
 D=R\oplus mX,\qquad X=x_5,\qquad m\in\ZZ_{\geq0}.
\]
Regularity of $R$ supplies the decomposition; finite multiplicity
and the absence of shifted copies use the common conformal vector.
The truncated level-thirteen fusion rule gives
\[
 X\boxtimes X=x_0\oplus x_1\oplus x_2\oplus x_3.
\]
In particular the odd summand cannot multiply into itself. If
$m>0$, the transformation $+1$ on $R$ and $-1$ on $mX$ preserves
every vertex product and has fixed algebra $R$. The compact-orbifold
tensor theorem applied to this order-two automorphism forces the
odd part to be an invertible simple $R$-module. This contradicts
the four displayed fusion channels (or
$\FPdim X=\sin(11\pi/15)/\sin(\pi/15)>1$). Therefore $m=0$.

For $V=B$ or $H$, positivity identifies
\[
 \ker_V(s_R)_1=\Com(\operatorname{Vir}(s_R),V).
\]
This kernel has conformal vector $\omega_R$, ambient integer grading,
and one-dimensional grade zero, so the extension obstruction gives
\[
 \Com(\operatorname{Vir}(s_R),V)=R,
 \qquad \Com(\Com(R,V),V)=R.
\]
Restriction gives the same result inside $U$. This argument does
not assume rationality of the opposite commutant. It isolates the
double-commutant input needed later for coset and mirror theorems.
The archived proof is
\artifact{hi_cft_20260923_next_structure_core_maximality.md}; its
external input is McRae's compact-orbifold tensor equivalence,
Corollary 4.8 in the source cited in that note.

\section{Residual conformal vectors and their complete finite algebra}
\label{expstructure:residual}

The residual stress splits as $s_R=q+r$ into commuting PCT-fixed
Virasoro vectors with
\[
 c_q=\frac{32}{19},\qquad c_r=\frac{182}{95},\qquad
 \frac{32}{19}+\frac{182}{95}=\frac{18}5.
\]
These central charges agree with those of $K(\mathfrak{sl}_2,17)$
and the diagonal affine coset at levels $(13,4;17)$, respectively.
Only the Virasoro fields, their membership and their commutation
are certified by
\artifact{hi_cft_20260923_proof_remaining_axes.json}; the larger
VOA identifications do not follow from these numbers.

The right raising and lowering currents have the exact simultaneous
weights
\[
 (h_{\mathcal T},h_Q,h_{s_R})
 =\left(\frac6{65},\frac7{65},\frac45\right).
\]
They sum to one. In particular their $R$ energy is $1/5$. This
small branching datum later identifies a tensor-generating core
label. The translation and PCT checks that justify using the actual
parafermion stress in this decomposition are stored in
\artifact{hi_cft_20260923_proof_remaining_coset.json}.

\subsection{Five ambient directions, but only two permitted in the proper algebra}

The full ambient grade-two kernel of $(\omega_R)_1$ in $B$ has
dimension five. Its producer enumerated a complete momentum ball
using rational LDL bounds and worked in the entire $1370$-dimensional
ambient grade, whose Cartan-zero block has dimension $86$. Thus this
is a full ambient result, rather than a kernel in $S_2$.
For $H$, the three lattice-glue sectors have grade-two dimensions
$1370,1377,1377$ and Cartan-zero dimensions $86,83,83$. Their
residual kernels have dimensions $5,4,4$, totaling $13$.

\begin{center}
\begin{tabular}{@{}lrrrr@{}}
\toprule
Ambient space&Grade two&Cartan zero&Residual kernel&Scope\\
\midrule
$B$&1370&86&5&Full grade\\
$H/B$ glue sector $1$&1377&83&4&Full grade\\
$H/B$ glue sector $2$&1377&83&4&Full grade\\
$H$&4124&252&13&Full grade\\
$S_2\subset U$&120&\text{selected block}&2&Generated span\\
\bottomrule
\end{tabular}
\end{center}
Sources:
\artifact{hi_cft_20260923_next_residual_ambient2.json} and
\artifact{hi_cft_20260923_next_residual_e8_ambient2.json}.
The complete momentum enumeration is part of the finite certification
boundary; neither count by itself identifies a VOA at every weight.

There is also a direct properness-sensitive test. In the full
five-dimensional residual space consider
\[
 \Phi(v)=\bigl(v_1E_+\bmod K_1,\ v_1E_-\bmod K_1\bigr).
\]
Its exact rank is three, and its kernel is exactly $\Span\{q,r\}$.
The kernel is separately checked to preserve all six known currents.
If $v\in C_U$ and $U_1=K_1$, then $v_1E_\pm\in U_1$, so
$\Phi(v)=0$. Since $q,r\in C_U$, this proves
\[
 (C_U)_2=\Span\{q,r\}
\]
once properness is supplied. This is a stronger statement than the
absence of three directions from the selected $S_2$: it excludes
them from the actual algebra using the current-preservation
condition. See
\artifact{hi_cft_20260923_next_structure_residual_current_kernel.json}.

The exact grade-three computation in the authenticated $S_3$ has
$58$ Cartan-zero dimensions, a residual kernel of dimension two,
and no residual Virasoro primary of weight three. Its kernel is
precisely $\Span\{Dq,Dr\}$. Using the separate conclusion $U_3=S_3$
therefore gives
\begin{equation}
 \ch C_U=q_{\mathrm{ser}}^{-3/20}
 \bigl(1+2q_{\mathrm{ser}}^2+2q_{\mathrm{ser}}^3
       +O(q_{\mathrm{ser}}^4)\bigr).
 \label{expstructure:residual-character}
\end{equation}
Here $q_{\mathrm{ser}}$ is the character variable, distinct from
the stress $q$. The source
\artifact{hi_cft_20260923_next_structure_residual_grade3.json} checks
the positive core modes as well as the zero-energy kernel. The
formula does not prove generation by $q,r$ or identify higher
character coefficients; any new residual strong generator must
start at weight at least four.

\subsection{The finite Griess algebra in a reproducible basis}
\label{expstructure:griess}

Use the exact basis $(q,r,v_1,v_2,v_3)$ of
\artifact{hi_cft_20260923_next_structure_residual_tensor.json}.
The last three basis vectors are obtained from the saved complete
ambient kernel by subtracting their invariant-form projections to
$q$ and $r$, in the producer's pivot order. This fixes their gauge;
the sparse lattice embedding of the five vectors remains in that
source. The following tables specify the entire finite algebra in
this basis without importing its mode engine.

Write $s=\sqrt{13}$, $t=\sqrt{-3}$ and abbreviate
\[
 [a,b,c,d]=a+bs+ct+dst.
\]
All four entries are rational. The invariant bilinear form is
$\langle x,y\rangle\one=x_3y$; it is not the positive Hermitian
Gram matrix in this non-real basis. In particular a complex or
negative-looking diagonal entry below is not a failure of vacuum
unitarity. We have
\[
\langle q,q\rangle=16/19,\qquad
\langle r,r\rangle=91/95,\qquad
\langle q,r\rangle=\langle q,v_i\rangle
 =\langle r,v_i\rangle=0.
\]
The residual weight-one space vanishes, so $x_2y=0$ for these
weight-two vectors and the Griess product $x\cdot y=x_1y$ is
commutative. Its unit is $(q+r)/2$. The producer checks every
product in the full ambient space and verifies invariance
$\langle x\cdot y,z\rangle=\langle x,y\cdot z\rangle$.

% The exact finite tensor tables are appended immediately below by a
% stdlib-only transcription of the small archived coefficient arrays.
\paragraph{Invariant form on the three complementary directions.}
\begin{longtable}{@{}ll@{}}
\toprule
Entry&Exact coefficient\\
\midrule\endhead
$\langle v_1,v_1\rangle$&$[-\frac{1350}{91},0,-\frac{1350}{91},0]$\\[3pt]
$\langle v_1,v_2\rangle$&$[\frac{470}{91},\frac{30}{91},\frac{30}{7},\frac{470}{91}]$\\[3pt]
$\langle v_1,v_3\rangle$&$[-\frac{820}{91},\frac{720}{91},-\frac{40}{91},-\frac{220}{91}]$\\[3pt]
$\langle v_2,v_2\rangle$&$[\frac{22940}{273},-\frac{8095}{546},-\frac{14145}{182},-\frac{170}{91}]$\\[3pt]
$\langle v_2,v_3\rangle$&$[-\frac{34960}{273},0,0,\frac{1840}{273}]$\\[3pt]
$\langle v_3,v_3\rangle$&$[\frac{22940}{273},\frac{8095}{546},\frac{14145}{182},-\frac{170}{91}]$\\[3pt]
\bottomrule
\end{longtable}
\paragraph{Action of the first stress.}
Write $q\cdot v_j=\sum_{i=1}^3Q_{ij}v_i$.
The row index is the output coordinate and the column index is the
input vector. Then $r\cdot v_j=2v_j-q\cdot v_j$, while
$q\cdot q=2q$, $r\cdot r=2r$ and $q\cdot r=0$.
\paragraph{The nine entries of $Q$.}
\begin{longtable}{@{}ll@{}}
\toprule
Entry&Exact coefficient\\
\midrule\endhead
$Q_{11}$&$\frac{4}{57}$\\[3pt]
$Q_{12}$&$[\frac{55}{1368},\frac{107}{1368},-\frac{29}{1368},\frac{35}{1368}]$\\[3pt]
$Q_{13}$&$[-\frac{71}{1368},\frac{1}{1368},\frac{13}{1368},-\frac{71}{1368}]$\\[3pt]
$Q_{21}$&$[-\frac{11}{190},-\frac{2}{95},\frac{1}{95},\frac{1}{190}]$\\[3pt]
$Q_{22}$&$[\frac{64}{57},-\frac{10}{57},0,0]$\\[3pt]
$Q_{23}$&$[-\frac{9}{76},\frac{5}{76},-\frac{9}{76},\frac{5}{76}]$\\[3pt]
$Q_{31}$&$[\frac{1}{76},-\frac{1}{380},\frac{13}{380},-\frac{1}{76}]$\\[3pt]
$Q_{32}$&$[-\frac{9}{76},-\frac{5}{76},\frac{9}{76},\frac{5}{76}]$\\[3pt]
$Q_{33}$&$[\frac{64}{57},\frac{10}{57},0,0]$\\[3pt]
\bottomrule
\end{longtable}
Its exact characteristic polynomial is
\[
 \det(xI-Q)=x^3-\frac{44}{19}x^2+\frac{449}{361}x-\frac{910}{6859}.
\]
The remaining six products are recorded coordinate by coordinate below.
Together with commutativity and the preceding stress actions, this is
the complete multiplication table. Define
$v_i\cdot v_j=c^q_{ij}q+c^r_{ij}r+\sum_k c^k_{ij}v_k$.
\paragraph{$v_1\cdot v_1$.}
\begin{longtable}{@{}ll@{}}
\toprule
Entry&Exact coefficient\\
\midrule\endhead
$c^{q}_{11}$&$[-5,0,-5,0]$\\[3pt]
$c^{r}_{11}$&$[-\frac{220100}{8281},0,-\frac{220100}{8281},0]$\\[3pt]
$c^{1}_{11}$&$[-\frac{409}{182},0,\frac{409}{182},0]$\\[3pt]
$c^{2}_{11}$&$[\frac{9}{91},\frac{9}{455},-\frac{9}{35},-\frac{9}{91}]$\\[3pt]
$c^{3}_{11}$&$[-\frac{198}{455},\frac{72}{455},-\frac{36}{455},\frac{18}{455}]$\\[3pt]
\bottomrule
\end{longtable}
\paragraph{$v_1\cdot v_2$.}
\begin{longtable}{@{}ll@{}}
\toprule
Entry&Exact coefficient\\
\midrule\endhead
$c^{q}_{12}$&$[\frac{475}{96},\frac{5}{96},\frac{65}{96},\frac{475}{96}]$\\[3pt]
$c^{r}_{12}$&$[\frac{319675}{49686},\frac{31925}{49686},\frac{31925}{3822},\frac{319675}{49686}]$\\[3pt]
$c^{1}_{12}$&$[\frac{2813}{1092},\frac{2117}{1092},-\frac{1919}{1092},-\frac{827}{1092}]$\\[3pt]
$c^{2}_{12}$&$[\frac{3}{14},\frac{38}{91},-\frac{3}{14},-\frac{38}{91}]$\\[3pt]
$c^{3}_{12}$&$[-\frac{1381}{728},-\frac{13}{56},-\frac{1381}{728},-\frac{13}{56}]$\\[3pt]
\bottomrule
\end{longtable}
\paragraph{$v_1\cdot v_3$.}
\begin{longtable}{@{}ll@{}}
\toprule
Entry&Exact coefficient\\
\midrule\endhead
$c^{q}_{13}$&$[-\frac{335}{96},\frac{715}{96},-\frac{205}{96},-\frac{235}{96}]$\\[3pt]
$c^{r}_{13}$&$[-\frac{782375}{49686},\frac{495475}{49686},\frac{47675}{49686},-\frac{143875}{49686}]$\\[3pt]
$c^{1}_{13}$&$[\frac{2813}{1092},-\frac{2117}{1092},\frac{1919}{1092},-\frac{827}{1092}]$\\[3pt]
$c^{2}_{13}$&$[\frac{1381}{364},-\frac{13}{28},0,0]$\\[3pt]
$c^{3}_{13}$&$[\frac{3}{14},-\frac{38}{91},-\frac{3}{14},\frac{38}{91}]$\\[3pt]
\bottomrule
\end{longtable}
\paragraph{$v_2\cdot v_2$.}
\begin{longtable}{@{}ll@{}}
\toprule
Entry&Exact coefficient\\
\midrule\endhead
$c^{q}_{22}$&$[\frac{10405}{72},-\frac{4265}{144},-\frac{5615}{48},\frac{85}{24}]$\\[3pt]
$c^{r}_{22}$&$[\frac{3607250}{74529},-\frac{366500}{74529},-\frac{1476500}{24843},-\frac{174250}{24843}]$\\[3pt]
$c^{1}_{22}$&$[-\frac{286577}{13104},-\frac{817}{1008},\frac{559}{1008},\frac{15083}{13104}]$\\[3pt]
$c^{2}_{22}$&$[-\frac{6779}{1092},-\frac{601}{1092},\frac{869}{364},\frac{21}{52}]$\\[3pt]
$c^{3}_{22}$&$[\frac{547}{364},\frac{41}{182},\frac{41}{14},\frac{547}{364}]$\\[3pt]
\bottomrule
\end{longtable}
\paragraph{$v_2\cdot v_3$.}
\begin{longtable}{@{}ll@{}}
\toprule
Entry&Exact coefficient\\
\midrule\endhead
$c^{q}_{23}$&$[-\frac{20425}{96},0,0,\frac{1075}{96}]$\\[3pt]
$c^{r}_{23}$&$[-\frac{1330475}{16562},0,0,\frac{70025}{16562}]$\\[3pt]
$c^{1}_{23}$&$[\frac{288059}{26208},\frac{15161}{8736},\frac{288059}{26208},-\frac{15161}{26208}]$\\[3pt]
$c^{2}_{23}$&$[\frac{4673}{1092},-\frac{11}{14},\frac{499}{273},-\frac{457}{1092}]$\\[3pt]
$c^{3}_{23}$&$[\frac{4673}{1092},\frac{11}{14},-\frac{499}{273},-\frac{457}{1092}]$\\[3pt]
\bottomrule
\end{longtable}
\paragraph{$v_3\cdot v_3$.}
\begin{longtable}{@{}ll@{}}
\toprule
Entry&Exact coefficient\\
\midrule\endhead
$c^{q}_{33}$&$[\frac{10405}{72},\frac{4265}{144},\frac{5615}{48},\frac{85}{24}]$\\[3pt]
$c^{r}_{33}$&$[\frac{3607250}{74529},\frac{366500}{74529},\frac{1476500}{24843},-\frac{174250}{24843}]$\\[3pt]
$c^{1}_{33}$&$[\frac{22027}{1872},-\frac{27935}{13104},-\frac{139655}{13104},-\frac{2231}{13104}]$\\[3pt]
$c^{2}_{33}$&$[\frac{547}{364},-\frac{41}{182},-\frac{41}{14},\frac{547}{364}]$\\[3pt]
$c^{3}_{33}$&$[-\frac{6779}{1092},\frac{601}{1092},-\frac{869}{364},\frac{21}{52}]$\\[3pt]
\bottomrule
\end{longtable}

\subsection{The residual charge spectrum and a unique charge-four-fifths axis}
\label{expstructure:axes}

The exact idempotent equations for
$e=aq+br+t_1v_1+t_2v_2+t_3v_3$ are $e\cdot e=2e$.
If its central charge is $c_e$, invariance against the total stress
makes
\[
 \frac{c_e}{2}=\frac{16}{19}a+\frac{91}{95}b,
 \qquad b=\frac{95c_e}{182}-\frac{80a}{91}.
\]
This linear elimination is useful for independently reconstructing
the polynomial system from the printed multiplication table. The
all-charge calculation found a zero-dimensional quotient of length
$32$. Multiplication by $c_e$ on that quotient has the factored
characteristic polynomial encoded by the following table. A pair in
the first two columns contributes both linear factors with the listed
multiplicity.

\begin{center}
\begin{tabular}{@{}ccc@{}}
\toprule
$c_e$&$18/5-c_e$&Multiplicity of each factor\\
\midrule
$0$&$18/5$&1\\
$4/5$&$14/5$&1\\
$6/7$&$96/35$&1\\
$52/55$&$146/55$&1\\
$34/35$&$92/35$&1\\
$6/5$&$12/5$&2\\
$10/7$&$76/35$&2\\
$58/35$&$68/35$&1\\
$32/19$&$182/95$&2\\
$96/55$&$102/55$&1\\
$62/35$&$64/35$&1\\
\midrule
\multicolumn{2}{c}{$c_e^2-(18/5)c_e+32168/10225$}&2\\
\bottomrule
\end{tabular}
\end{center}
The sum of the factor degrees with multiplicity is $32$.
Complementation $e\mapsto q+r-e$ explains the symmetry about
$9/5$. These are algebraic central charges of the complex idempotent
scheme. A real charge is not a certificate of a PCT-fixed stress,
and a Virasoro stress alone is not a rational conformal frame.
Scheme multiplicity must not be read as a count of distinct unitary
axes. The record is
\artifact{hi_cft_20260923_next_structure_residual_all_charges_v2.json}.

At $c_e=4/5$, the reduced exact equations give a unique point:
\begin{align*}
 e={}&\frac1{24}q+\frac8{21}r
      +\frac7{120}(1+t)v_1\\
 &+\frac{19+s+7t+st}{400}v_2
      +\frac{19-s-7t+st}{400}v_3.
\end{align*}
The original ideal computation, in
\artifact{hi_cft_20260923_next_structure_residual_potts.json},
only established this algebraic point. The separate verifier
\artifact{hi_cft_20260923_next_structure_residual_potts_verify.json}
then checked its full ambient Virasoro products, PCT reality,
commutation with the core, and commutation with its complement.
Its $e_1$ spectrum on the five-dimensional residual algebra is
\[
 0^{\oplus3},\qquad (2/5)^{\oplus1},\qquad 2^{\oplus1}.
\]
The complement has central charge $14/5$. The verifier's conclusion
is the unitary minimal Virasoro subalgebra $L(4/5,0)$ in $C_B$;
it does not assert a weight-three Potts extension inside $C_B$.

The field $e$ has nonzero complementary coordinates and fails
membership in $S_2$. More decisively, the current-kernel argument
shows that $e\notin C_U$ in the proper case. In fact every
idempotent in $\Span\{q,r\}$ is one of $0,q,r,q+r$, because the
two stresses commute and each satisfies $x\cdot x=2x$. Thus
$C_U$ has no charge-$4/5$ conformal vector under the established
properness hypotheses. This does not conflict with the original
Potts algebra $P\subset U$: $P$ is part of the core, not a
subalgebra of this residual commutant.

The tensor producer also performed a complete test within the
rank-one lattice Ising ansatz supported by the residual momenta.
It checked $21$ momentum pairs with arbitrary coefficients
$u,v$ satisfying $uv=1/16$, found no such axis, and left no
underdetermined cases. This excludes that ansatz only. It is not
an exclusion of every Ising conformal vector by an unrestricted
idempotent classification.

\section{Supplementary all-weight identification of the ambient residual cosets}
\label{expstructure:neighbor}

\begin{reviewnote}{Scope of the supplementary coset route}
The following route is additional verification material from
\artifact{hi_cft_20260923_next_neighbor_independent_audit.md},
\artifact{hi_cft_20260923_next_structure_affine_pair.md},
\artifact{hi_cft_20260923_next_structure_residual_identification.md},
and \artifact{hi_cft_20260923_next_structure_residual_category.md}.
It is not needed to replace the main manuscript's positive-stress
regularity or finite closure arguments. Its conclusions use the
actual parafermion embeddings, actual simple-current sectors,
lattice gluing, and the external affine-embedding, coset-regularity
and mirror theorems specified below. The archived internal audits
are evidence of review, not independent formal certification of
these additional theorems or their application. We reproduce the
steps so they can be verified in this draft.
\end{reviewnote}

\subsection{The two-neighbor lattice preserves exactly the relevant perpendicular part}

Let $\Lambda=E_8$ be the actual ambient lattice and let
$\Pi=\Span_{\mathbb R}\{a_0,H_R\}$.
Because $\det L=39$ is squarefree, $L=\Lambda\cap\Pi$ is
saturated: the square of a nontrivial saturation index would divide
$39$. Put
\[
 \Lambda_0=\{x\in\Lambda:(x,t_0)\in2\ZZ\},\qquad
 \Lambda'=\Lambda_0\cup(t_0/2+\Lambda_0).
\]
The parity map is nontrivial since $(a_0,t_0)=3$.
Thus $[\Lambda:\Lambda_0]=2$. The vector $t_0/2$ has even norm
six, pairs integrally with $\Lambda_0$, and has double in
$\Lambda_0$. Hence $\Lambda'$ is even, integral and unimodular
of rank eight, and is again an $E_8$ lattice.

Let $N=\Lambda\cap\Pi^\perp$. Every vector of $N$ already lies
in $\Lambda_0$. Conversely, a new-coset vector $t_0/2+x$ cannot
be perpendicular to $a_0$: that would require
$(x,a_0)=-3/2$, contrary to integrality of $\Lambda$.
Therefore
\begin{equation}
 \Lambda\cap\Pi^\perp=\Lambda'\cap\Pi^\perp=N.
 \label{expstructure:perp-preserved}
\end{equation}
This equality needs the displayed integrality argument. It is not
a general property of neighboring lattices.

The plane lattice changes to
\[
 L'=\ZZ(2a_0)+\ZZ(t_0/2).
\]
Set $\alpha'=t_0/2$, $\gamma'=4a_0-t_0/2$, and
$\beta'=(\alpha'+\gamma')/2=2a_0$. Then
\[
 (\alpha')^2=6,\qquad (\gamma')^2=26,
 \qquad(\alpha',\gamma')=0,
 \qquad G_{(\alpha',\beta')}=
 \begin{pmatrix}6&3\\3&8\end{pmatrix}.
\]
Both $L$ and $L'$ have determinant $39$.

\subsection{The sign of the discriminant transport is substantive}

The natural old discriminant generator is
$\mu=H_R/39$. Its natural new partner is
\[
 \lambda'=\mu+a_0
 =\alpha'/3+3\gamma'/13,
 \qquad q_{L'}(\lambda')=40/39\equiv1/39\pmod1.
\]
If $\mu+\eta$, $\eta\in N^*$, is an old glue lift, then
$\lambda'+\eta=(\mu+\eta)+a_0$ has even $t_0$ pairing
$1+3=4$. Therefore the old class $n\mu$ and the new class
$n\lambda'$ have exactly the same $N$ discriminant component.
Unimodularity and the order $39$ account for every glue sector.
The expression $-\alpha'/3+3\gamma'/13$ is also a discriminant
generator, but corresponds to $14\lambda'$ modulo $L'$.
It is not the natural transport of $\mu$. The archived audit
explicitly corrects an earlier sign proposal at this point.

The actual $39$-sector description of the core is
\[
 R=\bigoplus_{n\bmod39}
     V_{L+n\mu}\otimes P_{-n}\otimes Q_{3n}.
\]
The symbols $P_{-n}\otimes Q_{3n}$ here denote the actual
multiplicity subspaces in the corresponding $N$-lattice sectors,
not formal modules with matching characters. The gluing preserves
these subspaces, so it transports the sum to
\[
 R'=\bigoplus_{n\bmod39}
     V_{L'+n\lambda'}\otimes P_{-n}\otimes Q_{3n}
     \subset V_{\Lambda'}.
\]
The multiplicity-space intertwining operators are unchanged; only
the torus factors are replaced by factors with the same charge
fusion. The old product closure excludes extra output multiplicity
spaces. Compatible lattice cocycles may be chosen to agree on
$N$; their remaining discrepancy is a coboundary and rescales the
nonzero intertwining operators. Positive lattice forms can likewise
be chosen to extend the original form on $N$. These are the
field-level ingredients of the transport.

\subsection{Actual affine currents in the neighbor}

The two identities
\[
 13\lambda'-\alpha'-6\beta'=\alpha'/3,
 \qquad
 9\lambda'-\alpha'-4\beta'=\gamma'/13
\]
exhibit the charged sectors supplying the affine currents. Their
parafermion charges are $(-1,0)$ and $(0,1)$, respectively.
The lattice and parafermion energies are
\[
\begin{array}{c|cc|c}
\text{factor}&h_{\rm lattice}&h_{\rm parafermion}&\text{sum}\\\hline
\alpha'&1/3&2/3&1\\
\gamma'&1/13&12/13&1
\end{array}
\]
with one-dimensional parafermion tops. They give actual commuting
$A_{1,3}$ and $A_{1,13}$ current algebras. Their stresses are
$\omega_{\alpha'}+p$ and $\omega_{\gamma'}+z$.
The primitive plane field $e^{\beta'}$ supplies the order-two
simple-current extension, so the standard lattice--parafermion
reconstruction gives
\[
 R'=(A_{1,3}\otimes A_{1,13})
       \oplus(J_3\otimes J_{13}),
 \qquad h_{J_3\otimes J_{13}}=3/4+13/4=4.
\]
The reconstruction uses the actual PF3 and PF13 fields already
identified earlier. Matching their central charges would not supply
these affine raising and lowering vectors.

Both old and new cores have plane stress plus $p+z$. A vector
commuting with either core has zero plane-Heisenberg energy and
therefore lies in the common $V_N$ of
\eqref{expstructure:perp-preserved}. Positivity then gives the
all-weight equality
\[
 \Com(R,V_\Lambda)=\Com(R',V_{\Lambda'})
 =\{v\in V_N:p_1v=z_1v=0\}.
\]
The order-two affine extension shares the affine-pair stress, so
it has the same commutant. This proof does not infer a full
commutant from the finite weight-two kernel.

\subsection{Why the affine pair has a unique relevant conjugacy class}

An index-three $\mathfrak{sl}_2$ in $E_8$ has Cartan characteristic
of norm six. In Euclidean $E_8$ coordinates the norm-six vectors
have three types, up to signs and permutations:
\[
 (2,1,1,0^5),\qquad (1^6,0^2),\qquad
 ((3/2)^2,(1/2)^6),
\]
with counts $1344$, $1792$, $3584$. The $D_8$ Weyl group is
transitive on each type; reflection in
$(1/2,\ldots,1/2)$ connects the types. Thus the characteristic
has one $E_8$ Weyl orbit. For fixed $h$, the centralizer group
acts with an open orbit on every possible raising operator:
$\operatorname{ad}(e):\mathfrak g_0\to\mathfrak g_2$ is
surjective by $\mathfrak{sl}_2$ representation theory. Two open
orbits in the vector space $\mathfrak g_2$ coincide. This
upgrades the Cartan-orbit assertion to conjugacy of the first
$\mathfrak{sl}_2$ embedding.

For its standard representative, use the conformal inclusion
$G_{2,1}\otimes F_{4,1}\subset E_{8,1}$. The finite branching
\[
 248=(14,1)\oplus(1,52)\oplus(7,26)
\]
and
\[
 14\downarrow A_{1,3}\times A_{1,1}
 =(3,1)\oplus(1,3)\oplus(4,2),
 \qquad
 7\downarrow A_{1,3}\times A_{1,1}=(3,1)\oplus(2,2)
\]
show that the current centralizer is exactly
$F_4\oplus A_1$, both at level one. The finite Lie centralizer
alone is not the required VOA assertion. To verify the latter,
list the simple lowest weights:
\[
\begin{array}{c|l}
A_{1,3}&0,3/20,2/5,3/4\\
A_{1,1}&0,1/4\\
F_{4,1}&0,3/5
\end{array}
\]
Only the vacuum triple and $(3,1,0)$, $(2,0,26)$, $(1,1,26)$
have integral sums, the last three at weight one. Only the
vacuum triple has zero $A_{1,3}$ energy. Regularity and the
shared conformal vector therefore imply
\[
 \Com(A_{1,3},V_{E_8})=F_{4,1}\otimes A_{1,1}
\]
at every weight, with no additional affine-commutant extension.

The second $A_{1,13}$ projects to the external $A_1$ with index
zero or one. The allowed $F_4$ indices quoted from the
Penkov--Zuckerman table in the archived affine-pair note are
\[
 1,2,3,4,6,8,9,10,11,12,28,35,36,60,156,
\]
and that source states uniqueness by index. There is no index
thirteen and there is a unique index twelve. Hence the second
factor must use indices $12+1$. The actual residual commutant is
therefore identified, under these external embedding inputs, as
\begin{equation}
 C_H\cong\Com(A_{1,13},F_{4,1}\otimes A_{1,1}).
 \label{expstructure:CH}
\end{equation}

\subsection{Transporting the order-three action distinguishes the \texorpdfstring{$B$}{B} coset}

Let $\chi:\Lambda\to\ZZ/3$ have kernel the $E_6\oplus A_2$
lattice. The vectors $a_0,t_0$ have charge zero. The extension
$\chi'$ to $\Lambda'$ is fixed by
$\chi'|_{\Lambda_0}=\chi$ and $\chi'(t_0/2)=0$; the relation
$2(t_0/2)=t_0$ is compatible. On $N$ the action is unchanged,
and on the reconstructed core it is trivial. Thus its action on
$C_H$ fixes precisely $C_B$.

The index-twelve $\mathfrak{sl}_2$ in $F_4$ has adjoint branching
\[
 52=6V_2\oplus4V_4\oplus2V_6.
\]
Its full triple centralizer has zero-dimensional Lie algebra;
the nilpotent centralizer has dimension $6+4+2=12$.
The weighted diagram is $0200$ (orbit $F_4(a_3)$), with
characteristic $h=(4,2,2,0)$ of norm $24$. The external orbit
and component-group tables cited in the archived residual
identification note give the finite triple centralizer $S_4$.
The actual order-three element is nontrivial because its fixed
weight-two space has dimension five inside a space of dimension
thirteen. All order-three elements of $S_4$ are conjugate.

A representative is the grading automorphism of
\[
 F_{4,1}\supset A_{2,1}\otimes A_{2,2}.
\]
The two nontrivial simple-current sectors have current tops
$(3,\overline6)$ and $(\overline3,6)$, of dimension $18$ each.
The principal $A_1$ factors in the two $A_2$ algebras have levels
four and eight. Consequently the same actual fixed residual is
\begin{equation}
 C_B\cong\Com(A_{1,13},
             A_{2,1}\otimes A_{2,2}\otimes A_{1,1}),
 \label{expstructure:CB}
\end{equation}
with diagonal levels $4+8+1$. It is the order-three fixed algebra
of \eqref{expstructure:CH}. The conjugacy of the grading action,
rather than a comparison of central charges, is the decisive step.

The finite count gives a convenient check on the identification.
The nineteen currents of the numerator in \eqref{expstructure:CB}
restrict as $3V_2\oplus2V_4$. Horizontal singlets at grade two
come from the separate factors with multiplicities $1,2,1$ and
from the three cross terms with multiplicities $2,1,1$, giving
eight. The diagonal affine vacuum removes one; the two
complementary $V_2$ current tops remove two descendant singlets.
Thus the commutant has $8-1-2=5$ states at weight two.
The charge-$4/5$ stress from the $A_{2,2}/A_{1,8}$ coset has
spectrum $0^3,(2/5)^1,2^1$, agreeing with the exact Griess result.

\subsection{Regularity and the full category of the ambient residual}

The displayed standard cosets contain a common conformal frame
\[
 L(4/5,0)\otimes
 \Com(A_{1,12},A_{1,4}\otimes A_{1,8})\otimes L(34/35,0).
\]
The charges are $4/5$, $64/35$, $34/35$, totaling $18/5$.
The last factor is the unitary minimal model $M(14,15)$. The
middle factor is the level-four diagonal $A_1$ coset; its
regularity is the specific theorem of Lin (Theorem 5.7 at $k=8$)
invoked in the archived note. The unitary finite extension of this
regular frame yields $C_2$-cofiniteness of the ambient residual.
Together with core maximality, the strong rationality of the
ambient lattice VOA, and quasiprimary compatibility of the core
stress, McRae's coset theorem supplies strong rationality. This
order of argument avoids assuming rationality of arbitrary cosets
or an unstated categorical-dimension condition on an arbitrary
extension.

There is more information than regularity. The right currents
have $R$ energy $1/5$ and charges $\pm2$. Among the fourteen core
simples, the two lowest-weight-$1/5$ tops are $y_3$ with charges
$\pm2$ and $y_6$ with charges $\pm4$. Hence the currents force
$y_3$ to occur in $B\downarrow_R$. Its fusion square is
\[
 y_3\boxtimes y_3
 =(1\oplus\tau)\boxtimes
 (0\oplus2\oplus4\oplus6\oplus8\oplus10\oplus12),
\]
so it tensor-generates every core label. Mutual commutants and
the mirror theorem then give fourteen distinct simple
multiplicity spaces $N_i$ and the actual diagonal branching
\[
 B=\bigoplus_{i\in\operatorname{Irr}(R)}M_i\otimes N_i.
\]
Their braided subcategory is $\Rep(R)^{\rm rev}$, with inverse
twists because $B$ is integer graded. This is a statement about
actual branching modules, not merely fourteen proposed characters.

Splitting off this nondegenerate subcategory by the modular
centralizer theorem, and using the local-module category of the
existing extension algebra $B$, yields
\begin{equation}
 \Rep(C_B)\simeq
 \mathrm{Fib}^{\rm rev}\boxtimes\mathrm{SO}(3)_{13}^{\rm rev}
 \boxtimes\mathcal P(E_6\oplus A_2).
 \label{expstructure:CB-category}
\end{equation}
The last factor has group $(\ZZ/3)^2$ and form
$(2a^2+b^2)/3$, so there are $14\cdot9=126$ simple residual
modules. The extension algebra is diagonal in the first two
mutually reversed factors; its dimension is the sum of squared
core dimensions, which is positive. Thus the local-module
argument does not omit a zero-dimension exception.

If $U$ contained all of $C_B$, its right currents would still
force the $y_3$ summand. Fusion closure of mirror labels would
force all fourteen summands, each of which has multiplicity one
in $B$. Hence $U=B$. Conversely $U=B$ certainly contains $C_B$.
This explains structurally why a proper $U$ must have a proper
residual $C_U\subsetneq C_B$. It supplies no further regularity
or character claim for $C_U$ itself.

\section{Complete finite top data behind the universal energy bound}
\label{expstructure:tops}

The main proof states the sharp bound $h\le17/10$. This section
prints the charge characters and all partial-energy traces used
there, so a reviewer need not recover them from the checking script.
The coefficients are the saved exact data in
\artifact{hi_cft_20260923_next_structure_order13_top_check.json},
with the preceding lattice--Potts minimization explained in
\artifact{hi_cft_20260923_next_structure_top_bound_sharp.md}.

\subsection{Every lattice--Potts top, including both minimizing charges}

For $\epsilon=1,\tau$ the core energy at right Cartan charge $m$ is
\[
 \frac{m^2}{156}+\frac{\mathbf1_{m\text{ odd}}}{4}
                  +p_\epsilon(m),
 \qquad m\equiv j\pmod{13},
\]
where $p_1(m)$ is $0$ or $2/3$ according as $3\mid m$ or not,
and $p_\tau(m)$ is $2/5$ or $1/15$ in the same two cases.
An odd $m$ carries a complete left affine doublet, and an even
$m$ a singlet. The minimum is attained within $-12\le m\le12$.
For a positive $j$, the table combines the conjugate labels $j,-j$;
its dimension is the dimension of that pair of tops. The absolute
charges listed can belong to different signs of $j$, but both
signs of every listed charge occur in the combined pair.

\begin{longtable}{@{}llllr@{}}
\toprule
Label&$|m|$&Left multiplicity&Potts weight&Pair dimension\\
\midrule\endhead
$(1,0)$&$0$&$1$&$0$&1\\
$(\tau,0)$&$0$&$1$&$2/5$&1\\
$(1,1)$&$1,12$&$2,1$&$2/3,0$&6\\
$(\tau,1)$&$1$&$2$&$1/15$&4\\
$(1,2)$&$2$&$1$&$2/3$&2\\
$(\tau,2)$&$2$&$1$&$1/15$&2\\
$(1,3)$&$3$&$2$&$0$&4\\
$(\tau,3)$&$3,10$&$2,1$&$2/5,1/15$&6\\
$(1,4)$&$4,9$&$1,2$&$2/3,0$&6\\
$(\tau,4)$&$4$&$1$&$1/15$&2\\
$(1,5)$&$5,8$&$2,1$&$2/3,2/3$&6\\
$(\tau,5)$&$5,8$&$2,1$&$1/15,1/15$&6\\
$(1,6)$&$6$&$1$&$0$&2\\
$(\tau,6)$&$6,7$&$1,2$&$2/5,1/15$&6\\
\bottomrule
\end{longtable}

As a worked minimization, for $(1,1)$ the charges $m=1,-12$
both have energy $12/13$:
\[
 \frac1{156}+\frac14+\frac23=\frac{12}{13},
 \qquad \frac{144}{156}=\frac{12}{13}.
\]
Their left multiplicities are two and one. For $(\tau,3)$,
$m=3,-10$ give respectively
\[
 \frac9{156}+\frac14+\frac25=\frac{46}{65},
 \qquad \frac{100}{156}+\frac1{15}=\frac{46}{65}.
\]
Omitting these second minimizers would remove charge ten and
invalidate the later right-$\mathfrak{sl}_2$ inequalities.
The analogous second minimizer at $(\tau,6)$ restores charge
seven.

\subsection{The fourteen core tops as full Cartan characters}

Let $\ell$ record the left Cartan and $z$ the right Cartan.
Write $\sigma=\ell+\ell^{-1}$ and $[m]=z^m+z^{-m}$ for
$m>0$. These are Laurent characters, unrelated to the square-root
normalization of the parafermion primary. The complete top characters
are:
\begin{longtable}{@{}llr@{}}
\toprule
Module&Top Cartan character&Dimension\\
\midrule\endhead
$x_0$&$1$&1\\
$x_1$&$1$&1\\
$x_2$&$1+\sigma[3]$&5\\
$x_3$&$1+[2]+\sigma[3]+[6]$&9\\
$x_4$&$[6]$&2\\
$x_5$&$\sigma[1]+[4]+[6]+\sigma[9]+[12]$&14\\
$x_6$&$[4]+\sigma[9]$&6\\
$y_0$&$1$&1\\
$y_1$&$1+\sigma[5]+[8]$&7\\
$y_2$&$1+[2]+\sigma[3]+\sigma[5]+[8]+[10]$&15\\
$y_3$&$[2]$&2\\
$y_4$&$\sigma[1]+[2]+[4]+[6]+\sigma[7]$&14\\
$y_5$&$\sigma[1]+[4]$&6\\
$y_6$&$[4]$&2\\
\bottomrule
\end{longtable}

These are complete $R$-tops. They need not individually be complete
right-$\mathfrak{sl}_2$ representations, because the right raising
and lowering currents are in $U$ but are not currents of $R$.
For instance $x_5$ has charge twelve but no charge ten. A genuine
$U$-module top is a nonnegative integral sum of such core tops
whose combined character must be a finite-dimensional representation
of both current Lie algebras.

This gives a direct way to see the two inequalities used in the
bound. The multiplicities of right charge twelve and ten are
$x_5$ and $y_2$, so the usual nonincreasing multiplicity string
within an $\mathfrak{sl}_2$ representation forces $y_2\ge x_5$.
In the left-doublet sector the right charge-nine multiplicity is
$x_5+x_6$, while charge seven occurs $y_4$ times. Dividing out
the common doublet gives $y_4\ge x_5+x_6$.

\subsection{All partial-energy traces, with the parafermion contribution retained}

For one core top write $d$ for dimension,
$L=\Tr\omega_{A_{1,1},1}$, $Q_H=\Tr H_R(0)^2$,
$P_0=\Tr p_1$, and $Z_0=\Tr z_1$. The traces are taken over
that entire core top. The relation
\[
 Z_0=h_Rd-L-Q_H/156-P_0
\]
checks each row from the common core stress. On a complete right
$\mathfrak{sl}_2$ representation the affine Sugawara trace is
$3Q_H/164$. Multiplying by its positive-stress coefficient gives
$83Q_H/6084$. Hence the contribution to the total trace is
\[
 W=L+\frac{83}{6084}Q_H+\frac{25}{6}P_0+\frac{25}{26}Z_0.
\]
The following table is a transcription, not a fresh evaluation of
the research checker.

\begin{longtable}{@{}lrrrrrrr@{}}
\toprule
Label&$h_R$&$d$&$L$&$Q_H$&$P_0$&$Z_0$&$4680W$\\
\midrule\endhead
$x_0$&$0$&$1$&$0$&$0$&$0$&$0$&$0$\\[3pt]
$x_1$&$\frac{2}{15}$&$1$&$0$&$0$&$0$&$\frac{2}{15}$&$600$\\[3pt]
$x_2$&$\frac{2}{5}$&$5$&$1$&$36$&$0$&$\frac{10}{13}$&$10440$\\[3pt]
$x_3$&$\frac{4}{5}$&$9$&$1$&$116$&$\frac{4}{3}$&$\frac{268}{65}$&$56640$\\[3pt]
$x_4$&$\frac{1}{3}$&$2$&$0$&$72$&$0$&$\frac{8}{39}$&$5520$\\[3pt]
$x_5$&$1$&$14$&$2$&$720$&$4$&$\frac{44}{13}$&$148560$\\[3pt]
$x_6$&$\frac{4}{5}$&$6$&$1$&$356$&$\frac{4}{3}$&$\frac{12}{65}$&$54240$\\[3pt]
$y_0$&$\frac{2}{5}$&$1$&$0$&$0$&$\frac{2}{5}$&$0$&$7800$\\[3pt]
$y_1$&$\frac{8}{15}$&$7$&$1$&$228$&$\frac{4}{5}$&$\frac{92}{195}$&$36960$\\[3pt]
$y_2$&$\frac{4}{5}$&$15$&$2$&$472$&$\frac{8}{3}$&$\frac{56}{13}$&$110880$\\[3pt]
$y_3$&$\frac{1}{5}$&$2$&$0$&$8$&$\frac{2}{15}$&$\frac{14}{65}$&$4080$\\[3pt]
$y_4$&$\frac{11}{15}$&$14$&$2$&$312$&$\frac{8}{5}$&$\frac{14}{3}$&$81480$\\[3pt]
$y_5$&$\frac{2}{5}$&$6$&$1$&$36$&$\frac{2}{5}$&$\frac{10}{13}$&$18240$\\[3pt]
$y_6$&$\frac{1}{5}$&$2$&$0$&$32$&$\frac{2}{15}$&$\frac{4}{65}$&$4920$\\[3pt]
\bottomrule
\end{longtable}

For example the $x_5$ character gives
$Q_H=2(2\cdot1^2+4^2+6^2+2\cdot9^2+12^2)=720$.
Its four odd-charge doublets contribute $L=2$, while the Potts
trace is $4$. Therefore
$Z_0=14-2-720/156-4=44/13$, agreeing with the table.
This row alone has $W/d>17/10$; it is the compatibility
inequalities between complete core tops that lower the universal
bound. Bounding every individual partial energy separately would
lose that information.

\subsection{What makes the top bound an all-module statement}

The top of an arbitrary simple $U$-module is not assumed to have a
positive Hermitian form at this point. Instead regularity of $R$
as a weak-module theory decomposes the restriction into ordinary
$R$-modules. Intrinsic and ambient weights agree for vectors of
$R$, so every $R$ mode lowering its intrinsic degree also lowers
ambient degree and annihilates the lowest $U$ space. That space is
stable under all $R$ Zhu zero modes and is therefore a direct sum
of complete simple $R$-tops. Finite-dimensional current Lie algebra
representations are semisimple algebraically. These are the precise
inputs behind the character inequalities.

Writing $D_E$ for the total dimension and $W_E$ for the sum of
the listed $W$ contributions, the printed dual identity in the
main proof has only the nonnegative differences
$y_2-x_5$, $y_4-x_5-x_6$ and nonnegative core multiplicities.
It gives $W_E\le(17/10)D_E$. The stress identity gives
$W_E=hD_E$, so every simple $U$-module satisfies $h<2$.
No simultaneous diagonalization of all partial stresses, and no
unitarity of the unknown module, enters this trace argument.

The resulting absence of hidden branching summands is particularly
concrete. Every simple summand of an integer-graded lattice module
has integral lowest weight. Only zero and one are below two.
The vacuum multiplicity is one; the remaining $80$ currents form
one exact irreducible top. Thus $B=U\oplus X$ after properness.
For a mixed lattice module starting at weight one with an
irreducible $81$-dimensional top, an additional summand could only
start at an integer weight at least two, which the same bound
excludes. The finite irreducibility checks and the universal top
bound have different jobs: one identifies the lowest action,
the other rules out all later summands.

\section{The supplementary charge-six lattice coset and its limits}
\label{expstructure:c6}

The charge-six construction uses the same saturated plane lattice
$L$ but removes only its Heisenberg/lattice sector. Put
\[
 C_6=\Com(V_L,U),\qquad N=E_8\cap\Pi^\perp.
\]
Lattice positivity gives
$C_6=U\cap V_N=\Com(\mathcal H_\Pi,U)$, where
$\mathcal H_\Pi$ is the plane Heisenberg algebra. Its conformal
vector is the perpendicular stress. The zero-perpendicular-energy
space in $V_{E_8}$ is exactly $V_L$: a vector of zero energy has
neither perpendicular momentum nor perpendicular oscillators.
This proves $\Com(C_6,U)=V_L$ at every weight.

The source route is
\artifact{hi_cft_20260923_next_c2_asymmetric_coset.md}, with the
independent local audit
\artifact{hi_cft_20260923_next_structure_c6_coset_audit.md}.
Its use of the lattice-coset regularity theorem is specific: it
checks the positive-definite lattice double commutant required by
Section 4.3 of the Creutzig--Kanade--Linshaw--Ridout source before
applying its Lemma 4.6, Theorem 4.12 and Corollary 4.13.
It does not claim rationality for arbitrary cosets. Under the
strong rationality of $U$ already obtained in the main proof,
this route yields a strongly rational $C_6$. As elsewhere in
this supplementary route, these theorem applications are written
out for review and are not a new external validation.

\subsection{A fully specified rank-six lattice, not only a determinant}

The lattice $N$ has determinant $39$ and root system
$A_2\oplus A_2\oplus A_1$. Its full isometry class is the even
index-six extension
\[
 N=(A_2\oplus A_2\oplus A_1\oplus\langle78\rangle)
       +\ZZ(\omega,\omega,\alpha/2,u/6),
\]
where $\omega^2=2/3$, $\alpha^2=2$ and $u^2=78$. The glue
vector has norm $2/3+2/3+1/2+78/36=4$. In a basis
$(a_1,a_2,b_1,b_2,c,v)$ consisting of simple roots and this glue
vector, the Gram matrix is
\begin{equation}
\begin{gathered}
G_N=\begin{pmatrix}
2&-1&0&0&0&1\\
-1&2&0&0&0&0\\
0&0&2&-1&0&1\\
0&0&-1&2&0&0\\
0&0&0&0&2&1\\
1&0&1&0&1&4
\end{pmatrix},\\
\nu=6v-4a_1-2a_2-4b_1-2b_2-3c.
\end{gathered}
\label{expstructure:N-gram}
\end{equation}
with $\nu^2=78$ and $\nu$ orthogonal to the five roots.
The determinant check is the Schur complement
\[
 18\left(4-\frac23-\frac23-\frac12\right)=39.
\]
The letter $\nu$ here is the perpendicular vector denoted $u$
in the abstract lattice description, not the original weight-two
generating field.

There is also a basis in the original ambient coordinates.
Let $e_0,\ldots,e_7$ be the simple-root basis, with diagonal Gram
entries two and edges
$(0,1),(1,2),(2,3),(3,4),(2,5),(6,7)$ of pairing $-1$.
The index-three glue vector is
\[
 \lambda=(4/3,5/3,2,4/3,2/3,1,2/3,1/3).
\]
The vectors $e_0,\ldots,e_6,\lambda$ form an integral basis of
$E_8$. For $x=\sum_{i=0}^6y_ie_i+k\lambda$, orthogonality to
the plane says
\[
 y_5=-k,\qquad y_6=-\sum_{i=0}^4y_i-8k.
\]
Thus $n_i=e_i-e_6$, $0\le i\le4$, and
$n_5=\lambda-e_5-8e_6$ form a complete integral basis of $N$.
After replacing $n_5$ by
\[
 w=n_5-2n_0-2n_1-2n_2-n_3-n_4
   =(-2/3,-1/3,0,1/3,-1/3,0,2/3,1/3),
\]
the Gram matrix is
\[
\begin{pmatrix}
4&1&2&2&2&-2\\
1&4&1&2&2&-1\\
2&1&4&1&2&-1\\
2&2&1&4&1&0\\
2&2&2&1&4&-2\\
-2&-1&-1&0&-2&2
\end{pmatrix}.
\]
The archived proof identifies the isometry class by the root and
glue argument below; it does not claim to have saved a change-of-basis
matrix from this ambient basis to \eqref{expstructure:N-gram}.
All of this data is in
\artifact{hi_cft_20260923_next_c6_lattice.md}.

\subsection{The root-and-glue proof, including the excluded alternative}

The authenticated neighbor adjoint table has $22$ joint
Cartan-zero vectors. Removing the eight Cartan directions leaves
$14$ roots in $N$. No such root lies in the $E_6\oplus A_2$
root lattice: orthogonality to $a_0$ removes the $e_5$ coordinate
of an $E_6$ root, leaving the parabolic $A_5$; $H_R$ has positive
pairing two with every simple root of that $A_5$, and also with
both roots of the separate $A_2$. Thus all perpendicular roots
carry nonzero charge in $E_8/(E_6\oplus A_2)=\ZZ/3$.

An $A_3$ subsystem would have consecutive simple roots
$r,s,t$. Nonzero charges of $r+s$ and $s+t$ force the three
simple charges to agree; then $r+s+t$ has zero charge, a
contradiction. Hence every irreducible root component is $A_1$
or $A_2$. Fourteen roots and rank at most six leave
$A_2\oplus A_1^4$ or $A_2^2\oplus A_1$. The first has full
rank and determinant $48$, whose index in $N$ would have square
$48/39$, impossible. Therefore the rank-five root lattice is
$R_0=A_2^2\oplus A_1$.

Its discriminant form
\[
 D_{R_0}=(\ZZ/3)^2\oplus\ZZ/2,
 \qquad q(a,b,c)=\frac{a^2+b^2}{3}+\frac{c^2}{4}
\]
has no nonzero isotropic element, so $R_0$ is primitive in its
span. Let its rank-one perpendicular have norm $2m$, and let
$r$ be the glue index. Injectivity of the glue projections into
the two discriminant groups gives $r\mid6$. Determinants give
\[
 39=\frac{18(2m)}{r^2},\qquad m=\frac{13r^2}{12}.
\]
Integrality forces $r=6$ and $2m=78$. The order-six class on the
rank-one side has quadratic value $13/12\equiv1/12$; its partner
must have value $11/12$, forcing nonzero components in both
$A_2$ discriminants and in the $A_1$ discriminant. Diagram
conjugation reverses either $A_2$ component, so the glue is unique
up to isometry. This recovers the complete lattice above.

Its VOA has weight-one dimension $6+14=20$, with Lie algebra
$A_2\oplus A_2\oplus A_1$ and one abelian current. Its chiral
category is the pointed category on $\ZZ/39$ with form
$-n^2/39$. In contrast $(C_6)_1=0$ in the proper branch: the
non-Cartan $U$ currents have nonzero plane charges and the two
Cartan currents fail the positive Heisenberg-mode condition.
The known $U_2=S_2$ conclusion gives $\dim(C_6)_2=7$.

\subsection{Category factorization and the transported extension algebra}

Every one of the $39$ lattice labels occurs in the actual core
$R$, hence in $U$. The mirror theorem for the regular mutual
pair $(V_L,C_6)$ identifies the branching modules with the full
reverse pointed category $\mathcal P_{39}^{\rm rev}$, where
$q_{\mathcal P_{39}}(n)=n^2/39$. Modular centralizer decomposition
and the existing diagonal extension give
\begin{equation}
 \Rep(C_6)\simeq\mathcal P_{39}^{\rm rev}\boxtimes\Rep(U).
 \label{expstructure:c6-category}
\end{equation}
The actual $U$ algebra over $V_L\otimes C_6$ is the diagonal
algebra in $\mathcal P_{39}\boxtimes\mathcal P_{39}^{\rm rev}$,
tensored with the unit of $\Rep(U)$. Its local-module category
is the second factor, explaining \eqref{expstructure:c6-category}.
Under the main Haagerup-center identification the rank is
$39\cdot12=468$. Its multiplicative central charge is consistent:
the reverse pointed factor has charge six modulo eight and a
Drinfeld center has trivial anomaly.

Let $A_U$ be the actual extension algebra of the original
$U\subset H$. Transitivity of extension algebras identifies the
ambient algebra over $V_L\otimes C_6$ as
$\Delta_{\mathcal P_{39}}\boxtimes A_U$, including its
multiplication. Taking the $V_L$ vacuum multiplicity yields the
actual extension $C_6\subset V_N$ with algebra
$\one_{\mathcal P_{39}^{\rm rev}}\boxtimes A_U$.
Thus the $C_6$-preserving algebra-bimodule category is
\[
 \mathcal P_{39}^{\rm rev}\boxtimes
 {}_{A_U}\Rep(U)_{A_U}
 \simeq\mathcal P_{39}^{\rm rev}\boxtimes
 \mathcal F\boxtimes\mathcal F^{\rm op},
 \qquad \mathcal F=\Rep(U)_{A_U}.
\]
This transports the actual associativity data of the original
extension, not just its fusion matrices. The categorical defect
statement retains the full-CFT realization requirements: this
supplementary coset argument does not independently prove complete
unitarity of $C_6$ or the needed comparison of every bimodule with
a physical defect operator. Nor does it select $\mathcal F=\HH_3$
when the original extension is only known to have
$\mathcal F\simeq\HH_2$ or $\HH_3$.

\subsection{Precise obstructions to obtaining charge two by the immediate routes}

Several negative results constrain this additional construction.
They do not identify the physical chain or rule out an unrelated
charge-two theory.

First, the simple dimensions of $\Rep(U)$ are
$1,d^2,1+d^2$ (four times), and $3d$ (six times), where
$d=(3+\sqrt{13})/2$. Its only invertible is the unit and
$(3d)^2>1+d^2$. If two nontrivial fusion factors existed, each
would have a noninvertible simple of dimension at least $3d$;
their external product would exceed the largest allowed simple
dimension. Thus this category has no nontrivial fusion-category
tensor factorization. In particular a tensor factorization of
$U$ into unitary strongly rational charges two and six would
force one factor to be holomorphic, whose charge must be zero
modulo eight, a contradiction. This does not exclude a finite
conformal extension of a tensor product or a mutual-commutant pair.

Second, the transparent current algebra inside $U$ is contained
in its two $\mathfrak{sl}_2$ factors. Their full Sugawara charge is
\[
 1+\frac{3\cdot39}{39+2}=\frac{158}{41}.
\]
A compact reductive current subalgebra has no larger Sugawara
charge, by positivity of the complementary affine coset.
An ordinary coset by such currents therefore has charge at least
$8-158/41=170/41>2$. The precise premise is that the gauged
currents commute with the specified symmetry action; generalized
quotients or modified actions are outside this statement.

Third, the exact charge-one and charge-two idempotent exclusions
in the complete seven-dimensional $(C_6)_2$ are recorded in
\artifact{hi_cft_20260923_next_c2_reduced_exclusion_c1_checked.json}
and
\artifact{hi_cft_20260923_next_c2_reduced_exclusion_c2_checked.json}.
Their original finite scope becomes exhaustive only after using
$U_2=S_2$ and the complete Cartan-commutant dimension. If $e$ is
any complex Virasoro vector in this weight-two space and $t_6$
is the total stress, the absence of weight-one states and
invariance give
\[
 (t_6,e)=(e,e)=c_e/2,
 \qquad t_6-e\text{ is a commuting Virasoro vector of charge }6-c_e.
\]
Thus the exclusions also rule out charges five and four. In
particular an ordinary conformal charge-four coset inside this
$C_6$ cannot leave charge two. The argument does not require
$e$ to be PCT real. The polynomial system and its finite
certificate interpretation are described in
\artifact{hi_cft_20260923_next_c2_slice_theorem.md}.

These facts distinguish an actual charge-six lattice commutant
from an asserted charge-two physical theory. Finite topological
gauging itself preserves the stress and central charge. Deleting
a commuting conformal sector is additional structure and does not
automatically transport the desired defect action.

\section{Categorical exclusions: diagnostic calculations and their scope}
\label{expstructure:categorical-diagnostics}

The main branching argument eliminates the rank-four and higher
multiplicity rank-six possibilities before invoking unitary
classification. Several early notes only compared with familiar
affine models; the later notes give universal replacements.
This section records the small diagnostic calculations that are
particularly useful when checking that upgrade.

\subsection{The rank is an endomorphism-algebra calculation}

From the actual extension restriction
$H=\one\oplus X\oplus Y\oplus Y^*$, the alternatives are
\[
\begin{array}{c|c|c}
\text{case}&\End_U(H)&\text{boundary rank}\\\hline
Y\not\cong Y^*&\CC^4&4\\
Y\cong Y^*&\CC\oplus\CC\oplus M_2(\CC)&6
\end{array}
\]
The dimension is the sum of squared multiplicities, not the number
of distinct summands. In the second case the composition algebra
is noncommutative, an essential input to the rank-six fusion-ring
analysis. Center induction identifies this algebra with the
complexified fusion algebra of the all-module boundary category.
The local modules are a different category: for a holomorphic
extension they form $\Vec$ and cannot give this rank calculation.

The actual order-three lattice grading automorphisms inject into
the invertible objects of the boundary category. Untwistedness
requires additional information. The metric group of $B$ has
quadratic form $(2a^2+b^2)/3$; its two transverse isotropic lines
are $(a,a)$ and $(a,-a)$. Induction from the second line to modules
of the first-line extension is fully faithful by adjunction:
only the trivial intersection can contribute a morphism.
The target dimension is $9/3=3$, so the three images exhaust it.
This gives the tensor equivalence with untwisted
$\Vec_{\ZZ/3}$. The existence of three invertible objects alone
would not determine their associator.

\subsection{Why the balanced multiplicity-six near-group case is impossible}

For a rank-four $\ZZ/3+6$ near-group category, put
$d=3+2\sqrt3$ and $r=1+d$. In the balanced induction pattern,
the known central dimensions are
\[
 1,\quad 1+2d,\quad 2r\ (3\text{ times}),
 \quad r\ (4\text{ times}),\quad n_jd,
 \qquad \sum_jn_j^2=22.
\]
The normalized modular vacuum row has entries
\[
 b,\ a,\ \underbrace{1/3,1/3,1/3}_{3},
 \underbrace{1/6,1/6,1/6,1/6}_{4},
 \frac{n_j\sqrt3}{12},
 \qquad
 a=\frac{2+\sqrt3}{12},\quad b=\frac{2-\sqrt3}{12}.
\]
Signed Galois symmetry under $\sqrt3\mapsto-\sqrt3$ forces the
row of the unique dimension-$1+2d$ object to be
\[
 a,\ b,\ 1/3,1/3,1/3,\ 1/6,1/6,1/6,1/6,
 -n_j\sqrt3/12.
\]
That object is self-dual because its dimension is unique. Its
self-fusion coefficient is therefore
\[
 N_{XX}^{X}
 =\frac{a^3}{b}+\frac{b^3}{a}+\frac{3}{9}
        +\frac4{36}-\frac{22}{48}
 =\frac{97}{72}+\frac{32}{72}-\frac{33}{72}
 =\frac43.
\]
For a hand check, $a+b=1/3$, $ab=1/144$, so
$a^2+b^2=7/72$ and
$(a^3/b)+(b^3/a)=97/72$. Nonintegrality excludes the balanced
case without specifying any unknown twist or unprinted modular
row. The source is
\artifact{hi_cft_20260923_next_structure_m6_balanced_exclusion.md}.

For the remaining unbalanced case, the later universal projection
uses only a forced three-dimensional character subspace. With
$d=3+2\sqrt3$, its remaining twists $z_j$ satisfy
\[
 \sum_jn_j^2=6,\qquad
 \sum_jn_j^2z_j=-2-4i,\qquad
 \sum_jn_j^2z_j^2=-2+4i(1-\epsilon),
 \quad\epsilon\in\{1,-1\}.
\]
The second-moment modulus excludes $\epsilon=-1$. Next
\[
 \sum_jn_j^2(1+\Re z_j)(1+\Im z_j)
 =6-2-4+\tfrac12\Im(-2)=0.
\]
Each summand is nonnegative, so every twist is $-1$ or $-i$;
the total squared multiplicities are two and four respectively.
This argument allows a dimension-$2d$ object at the latter twist;
it does not assume exhaustion by a familiar rank-nine model.
In the categorical twist convention, before the common character
phase $e^{-2\pi i/3}$ is included, the stable weighted projection obeys
\[
 S_v=-\frac12\begin{pmatrix}1&2&1\\1&0&-1\\1&-2&1\end{pmatrix},
 \qquad T_v=\operatorname{diag}(1,-i,-1).
\]
The harmonic-theta coefficient calculation in the two orientations
then gives either $\dim U_1=8+A$, $A\in\ZZ_{\geq0}$, or
$\dim U_1=5+12P$, $P\in\ZZ_{\geq0}$. Neither can be six.
The universal scope is supplied by
\artifact{hi_cft_20260923_next_structure_m6_universal_projection.md};
the earlier
\artifact{hi_cft_20260923_next_structure_m6_character_obstruction.md}
by itself only excluded the two explicit ribbon models compared there.

\subsection{The exceptional rank-six multiplicity-two check}

For the noncommutative rank-six family,
$d^2=1+3md$ and the residual central objects have dimensions
$3n_jd$. The first two indicator moments are
\[
\begin{split}
 \sum_jn_j^2&=4m^2+2,\\
 \sum_jn_j^2\theta_j
   &=-\frac{m^2+m\sqrt{9m^2+4}}2,\\
 \sum_jn_j^2\theta_j^2
   &=\frac{-m^2+(2\epsilon-m)\sqrt{9m^2+4}}2.
\end{split}
\]
The root-count estimate eliminates $m>1$ except possibly the
squareclass ten. In that squareclass, a multiset summing to
$-a-b\sqrt{10}$ has at least $a+8b$ roots. Its equality statement
requires that the sum of squared roots belongs to $\QQ(\sqrt{10})$;
this is an actual hypothesis, provided here by the displayed
second indicator moment. Without it, Galois averaging would not
recover the original squared sum.

The last candidate is $m=2$, $\sqrt{9m^2+4}=2\sqrt{10}$.
It has $18$ roots, $a=2$, $b=2$, and attains the lower bound
$2+8\cdot2=18$. Equality forces squared sum $2$. The indicator
formula instead gives $-2$ when $\epsilon=1$ and
$-2-4\sqrt{10}$ when $\epsilon=-1$. This explicit contradiction
removes the exceptional case.

The corrected proof and its prior audit are
\artifact{hi_cft_20260923_next_structure_rank6_indicators.md}.
The same note records a useful failed diagnostic: the HLPW
primary-$n$ positivity tests are nonnegative for the entire family,
so those tests cannot prove $m=1$. A criterion that every candidate
passes is not evidence that it selects the desired multiplicity.

At $m=1$, the remaining six-root sum is
$(-1-\sqrt{13})/2$. The minimal quadratic Gauss orbit forces
six distinct thirteenth-root twists with exponents
$\{2,5,6,7,8,11\}/13$, each with $n_j=1$; the squared sum is
$(-1+\sqrt{13})/2$ and forces second indicator $+1$.
This gives the twelve central dimensions and twists. Determining
these modular data and proving a braided category equivalence
remain separate steps. The main manuscript's unitary route uses
the separately stated unitary classification; the older optional
pivotal route invokes a broader external exhaustive census whose
equation-solving files were not independently replayed here.

\section{Analytic verification map: domains, sectors and extensions}
\label{expstructure:analytic-map}

The analytic part of the current manuscript proves three bridges:
A2 is the independent net sector count; A3 reconstructs and exhausts
ordinary unitary modules; A1 compares the particular bounded annulus
insertions. Their roles differ. The following table records both
the conclusion and the premise that would be lost by a shortcut.
The references point to the proofs already present in this draft.

\begin{longtable}{@{}p{0.19\linewidth}p{0.36\linewidth}p{0.37\linewidth}@{}}
\toprule
Step&Actual conclusion&Required distinction\\
\midrule\endhead
Vacuum form&The PCT-stable $U$ has a positive invariant form.&This
is a form on the vacuum VOA, not yet on every simple module.\\
Finite index&The bounded-insertion subnet is completely rational;
its abstract vacuum isomorphism transfers that property to
$\mathcal N_U$.&An isomorphism-invariant property does not identify
individual concrete insertion operators.\\
A2&The actual $H$ restriction has multiplicities $1,1,2$; the
boundary ring gives twelve net sectors.&The count is derived for
the net category, rather than copied from the VOA modular matrix.\\
A3&Each irreducible net sector reconstructs an irreducible ordinary
unitary $U$-module; the functor is fully faithful and linear $*$.&No
tensorator or braiding compatibility is claimed for this functor.\\
Exhaustion&Twelve distinct net images exhaust twelve algebraic
simple modules.&Equality of counts gives linear surjectivity,
not positivity of canonical fusion pairings.\\
A1&The two concrete nets coincide and each localized annulus
insertion transports to its bounded module form.&Central path lifts,
range domains and full insertion interiors are part of the claim.\\
Fusion positivity&Tener's theorem gives positivity of canonical
VOA fusion forms after strong unitarity and A1.&Unitarizability of
the fusion product is supplied before applying that theorem.\\
Classification&The now unitary all-module category has the
specified Haagerup classes and the center is $Z(\HH_3)$.&Matching
modular data alone is not this tensor classification.\\
\bottomrule
\end{longtable}

\subsection{Closed-field reduction: the domain fact needed before irreducibility}

For a closable smeared field $T_0$ on a common finite-energy core
$D$ and an orthogonal projection $P$, the two statements
$PD\subset D$ and $T_0P=PT_0$ on $D$ imply that $P$ reduces
the closure and its adjoint. A graph-norm approximation establishes
$TP=PT$ on the closed domain; applying the same argument to
$1-P$ decomposes the closed operator as a direct sum. Its polar
partial isometry and the spectral projections of its absolute
value are then block diagonal. This is the reason the formal
module projection becomes an intertwiner of the generated local
von Neumann algebras. Formal mode commutation without domain
invariance would not give that conclusion.

In reverse reconstruction, choose one ordinary simple algebraic
summand $S$ of the reconstructed weak module $M$ and let $P$
project onto its Hilbert closure. The needed domain inclusion is
\[
 P1_{[0,R]}(L_0)\mathcal K
   \subset\bigoplus_{h+n\le R}S_{h+n}\subset S.
\]
It uses finite-dimensional energy spaces of $S$, not a prior
finite-multiplicity assertion for the entire sector Hilbert space.
The invariant form then gives mode commutation. A smeared field
need not map $M$ into itself; its Hilbert-space value and boundedness
of $P$ suffice to pass to the closed field. Irreducibility forces
$P=1$ and hence $M=S$. These points are the substantive content
of Proposition~\ref{clm:analytic-reconstruction}.

Full faithfulness also needs boundedness of an algebraic module
isomorphism. Compare its two positive invariant forms grade by
grade. The resulting positive grade-preserving comparison operator
is a $U$-endomorphism. An eigenvector in the finite lowest space
generates an eigenspace submodule, so simplicity makes the operator
one scalar on the whole module. The isomorphism therefore has a
bounded unitary rescaling. This is stronger than simply observing
that algebraic Schur's lemma gives one-dimensional morphism spaces.

\subsection{A fractional-energy test for the common-lift argument}

Vacuum normalization cannot distinguish every universal rotation
lift: $Z=e^{2\pi iL_0}$ is the identity on the integer-graded
vacuum but acts by $e^{2\pi ih}$ on a simple module of lowest
weight $h$. For example, on a module with $h=2/13$ it acts by
$e^{4\pi i/13}$. Identifying two annulus lifts only from their
vacuum scalar would miss this phase.

The pairwise construction instead first chooses the lifts of the
two localized annuli in the canonical localized extension, whose
remaining ambiguity is scalar. It normalizes their nonzero vacuum
operators there, and then undoes the same universal lift
$\widetilde\delta$ of a common localizer. On a module, changing
that common lift by $Z^m$ changes the pair as
\[
 A_M\longmapsto Z^{-m}A_M,
 \qquad B_M\longmapsto B_MZ^m.
\]
Every $U$ mode changes conformal energy by an integer, so it
commutes with $Z$. Hence
\[
 B_MZ^mY_M(v,z)Z^{-m}A_M=B_MY_M(v,z)A_M.
\]
This is a pairwise cancellation at fractional as well as integer
lowest weights. It neither assumes a global choice of lifts for
all individual annuli nor discards a deck transformation by
vacuum normalization.

For the semigroup part of an incoming annulus, the framing
$h(\theta,s)=\varphi_{t-s}(e^{i\theta})$ has logarithmic velocity
$i\rho(e^{i\theta})\partial_\theta$. With the stated stress
convention its integrated action is exactly $e^{-tL(\rho)}$.
Thin annuli supply the boundary reparametrizations, giving the
specific operator
$U(\psi)e^{-tL(\rho)}U(\gamma)^*$, including its scalar.
This concrete identification is what an abstract vacuum-net
isomorphism alone does not provide.

\subsection{Range, density and holomorphic continuation are three separate checks}

For a geometrically complementary annulus $C_M$, round
factorization gives
\[
 A_MC_M=\lambda r^{L_0},\qquad0<r<1,\quad\lambda\ne0.
\]
Thus every finite-energy vector lies in $\operatorname{Ran}A_M$:
\[
 v=A_MC_M(\lambda^{-1}r^{-L_0}v).
\]
Injectivity of $A_M$ gives the equality of test domains
\[
 A_M^{-1}M^{\rm fin}=C_MM^{\rm fin}.
\]
Boundedness and dense range of $C_M$ then give density of this
domain. The outgoing adjoint has the same properties. These are
range and density statements; neither follows merely from saying
that an annulus operator is bounded.

On those two test domains, both the input $A_M\xi$ and output
$B_M^*\eta$ have finite energy support. For a weight-$k$ state
$v$, only modes with $q=p+k-n-1$ contribute between input energy
$p$ and output energy $q$. Hence
\[
 \langle Y_M(v,z)A_M\xi,B_M^*\eta\rangle
\]
is a finite Laurent polynomial in $z$. Sewing first identifies
it with the bounded geometric insertion on the incoming interior;
adjoints do so on the outgoing interior. The holomorphic identity
theorem extends it to the interior of the union of the two closed
regions. Every component meets one of the original open regions,
since the remaining set lies in their boundaries and has no open
component. Uniqueness on dense test domains identifies the bounded
operator, without treating the unsmeared formal field as everywhere
defined.

The prescribed insertion interior is the union over all compatible
geometric representatives of the same generalized operator pair.
The proof repeats over that union and uses uniqueness on overlaps.
Pairs with empty insertion interior still contribute the generator
$B_UA_U$, handled directly by local annulus transport. Finally,
locality and Haag duality on the same vacuum Hilbert space give
\[
 \mathcal A_U^{\rm bl}(I)\subset\mathcal N_U(I)
 \subset\mathcal A_U^{\rm bl}(I')'
 =\mathcal A_U^{\rm bl}(I).
\]
That is concrete equality of the two net constructions, the
conclusion needed for the module insertion representation.

\subsection{Primary theorem chain retained by the current A1 proof}

The following is the exact source map recorded by the final
27 September analytic revision. It is useful for reviewing the
pairwise argument without reconstructing its citation chain.

\begin{longtable}{@{}p{0.23\linewidth}p{0.68\linewidth}@{}}
\toprule
Source&Specific uses\\
\midrule\endhead
Tener, fields \cite{TenFields}&Definition 3.16 for incoming
semigroup data; Lemma 3.19 and Corollaries 3.20--3.22 for a common
localizer; Definition 4.2 for module range and density; Proposition
4.13 for vacuum covariance; proof of Theorem 4.28 for the union of
representative interiors; Section 5.1, formula (5.2), for paired
module annulus actions.\\
Henriques--Tener, annuli \cite{HTAnn}&Remark 4.2, formula (5.2)
and Theorem 5.5 for the integrated framing operator; Section 7
and Theorem 7.4 for the localized extension and local sector
transport.\\
Henriques--Tener, nets \cite{HTNet}&Lemmas 4.15 and 7.9 for
injectivity/dense range and round factorization; Proposition 4.17
for transport; Lemma 6.8 for local point insertions; formula (76)
for sewing; Corollary 6.15 for holomorphicity; Proposition 6.19 and
formula (81) for adjoints; Theorems 8.9 and 9.5(i) for vacuum
reconstruction and module insertions.\\
Tener, fusion \cite{TenFusion}&Definition 2.18 is the required
module representation property; Theorem 4.4 and Corollary 4.6
supply canonical fusion positivity after the separately proved
strong unitarity.\\
\bottomrule
\end{longtable}

The frozen earlier part of
\artifact{hi_cft_c8_revision_20260927_analytic.md} records an
unsuccessful search for a ready-made theorem proving A1. Its final
section explicitly supersedes that assessment with the pairwise
construction. Both records remain useful: the failed shortcut
explains why complete unitarity is not deduced from a linear
category equivalence, while the final proof supplies the additional
operator comparison. The current manuscript treats A1 as proved
under its ordinary construction inputs, not as an additional
unresolved assumption. External expert verification of the written
argument is still distinct from those internal review records.

\subsection{The three extension objects that must remain distinct}

\begin{longtable}{@{}p{0.2\linewidth}p{0.34\linewidth}p{0.37\linewidth}@{}}
\toprule
Object&Established construction&Identification boundary\\
\midrule\endhead
Original $H=V_{E_8}$&The given lattice inclusion over $U$, with
restriction $\one\oplus X\oplus Y^{\oplus2}$.&Its all-module
category is $\HH_2$ or $\HH_3$; the proof does not choose which.\\
Chosen net $\mathcal E_3$&Transport the canonical $\HH_3$ Q-system
to $\Rep(\mathcal N_U)$ and apply local net reconstruction.&It is
holomorphic of charge eight; no classification of arbitrary
holomorphic charge-eight nets as the $E_8$ lattice net is asserted.\\
Chosen VOA $H'$&After complete unitarity, transport the canonical
$\HH_3$ algebra to $\Rep(U)$ and use unitary VOA extension theory.&
Holomorphic charge-eight VOA classification identifies $H'$ with
$V_{E_8}$ as a VOA, but does not identify its embedding over $U$
with the original lattice embedding.\\
\bottomrule
\end{longtable}

For the chosen $\HH_3$ extension, all right algebra modules form
$\HH_3$, local modules form $\Vec$, and algebra bimodules form
$\HH_3\boxtimes\HH_3^{\rm op}$ with $36$ simples. Only the local
subcategory inherits the ambient braiding by descent. The
unitarization argument for all modules uses the braided left
action and the relative tensor product; it does not impose the
extra local-module monodromy condition. The diagonal full net
theory, by contrast, has partition function
$\sum_i|\chi_i|^2$ and uses the twelve-sector linear character
comparison. These constructions have different roles even though
they start from the same smaller chiral algebra.

\section{Structural provenance, replay results and historical status changes}
\label{expstructure:provenance}

A saved exact identity, an exact replay of its producer, and an
all-weight implication are different pieces of evidence. This
section records their relationship. No research producer was run
for the present expansion: the coefficients and results here were
read from the existing artifacts, with lightweight text and JSON
transcription only.

\subsection{The bounded structural replays of 27 September}

\begin{longtable}{@{}p{0.22\linewidth}p{0.47\linewidth}p{0.2\linewidth}@{}}
\toprule
Replay&Exact assertions recovered&Recorded resources\\
\midrule\endhead
Potts primary&Stress charge $4/5$; actual weight-three primary;
membership after a rank-$40$ descendant space and one additional
zero-mode pair, reaching rank $41$; $w_{P,2}$ eigenvalues $0,2,-2$
on $H_R,E_+,E_-$.&$8.234$ s; producer peak $97820$ KiB.\\
PF13 membership&The saved primary lies in the generated rank-$58$
Cartan-zero space after $154$ zero-weight pairs; PCT-stable primary
line; positive-stress coefficients $1,3403/4563,25/6,25/26$.&
$180.622$ s; adapter peak $120436$ KiB.\\
Order-thirteen core&$\dim S_2=120$; known tensor-core grade two
$15$; zero-residual-energy dimension $25$; ten extras and all six
one-dimensional witness charge spaces.&$6.455$ s; producer peak
$40216$ KiB.\\
Four mixed tops&All four sign choices complete; seed index $68$,
one-dimensional seed Cartan eigenspace, exact weighted adjoints,
cyclic dimension $81$, irreducible top.&$3.031$ s; producer peak
$26496$ KiB.\\
Top-bound arithmetic&All fourteen minima, dimensions, stress traces
and charge characters; both $\mathfrak{sl}_2$ inequalities and the
nonnegative dual identity.&Tiny exact Fraction calculation;
no field action is inferred.\\
\bottomrule
\end{longtable}

The files are, respectively,
\artifact{c8_revision_20260927/potts_replay.json},
\artifact{c8_revision_20260927/pf13_membership_replay.json},
\artifact{c8_revision_20260927/core_extension_replay.json},
\artifact{c8_revision_20260927/top81_replay.json}, and
\artifact{c8_revision_20260927/top_bound_replay.json}.
The first two corresponding logs end with the adapter status
\texttt{unchanged\_producer\_exact\_replay\_passed}. The adapter
compares complete mathematical outputs with the frozen certificate,
ignoring only run time and memory accounting. It authenticates
source and original certificate before and after, and redirects
exactly one requested output path and one write. This status is
stronger than an intermediate producer success message.

The core and top producers can record a bounded partial run, so
process exit zero alone is insufficient. The saved mixed-top record
has \texttt{status=complete}, four sectors, and in every sector
\texttt{cyclic\_dimension=81},
\texttt{seed\_eigenspace\_dimension=1},
\texttt{exact\_weighted\_adjoint\_checks=true}, and
\texttt{irreducible=true}. Even these full finite checks do not
supply the $h<2$ all-weight branching step by themselves.

The archived resource review
\artifact{c8_revision_20260927/structural_replay_review.md}
estimated below $0.5$ GiB for core and mixed-top jobs and below
$0.75$ GiB for the two membership jobs, including sparse pivots,
Fraction temporaries, copies and bounded mode caches. It used
serial root-controlled launches with a $2$ GiB address cap,
$8$ GiB process-group RSS watchdog, two CPUs/threads and the
permanent $32$ GiB host reserve. These measured historical jobs
do not authorize a larger future replay or renewal of an expired
resource exception.

\subsection{PF13 checkpoints and the boundary of the replay claim}

The full PF13 identification combines more than the membership
record. Its direct OPE stage used a continuation of a partially
completed run: fifteen completed exact assertions were retained
in an authenticated log, and the remaining direct coefficient was
computed by a separate continuation. The final simple pole is
obtained by a checked algebraic Jacobi/skew-symmetry derivation,
not by treating an interrupted weight-seven expansion as zero.
The weight-six decoupling identity had already completed before
that direct high-weight cross-check was stopped.

The reconstruction therefore depends on the recursive field
checkpoints
\artifact{hi_cft_20260923_proof_pf13_W5.json} and
\artifact{hi_cft_20260923_proof_pf13_W6.json}, as well as the
exact-field, DLY-conversion and weight-seven producers. A source
provenance audit found these five dependencies missing from an
earlier combined graph. Their historical hashes were retained in
\artifact{hi_cft_20260923_proof_artifact_audit.json}; the revised
\artifact{c8_revision_20260927/authenticate_inputs.py} explicitly
binds them. The $W^6$ file is about $8.8$ MiB and is stream-hashed
rather than parsed by the revision's $8$ MiB JSON cap.

Hash authentication proves identity of bytes with the recorded
checkpoint, not correctness of the modes that produced it. The
membership replays above do not rerun the expensive $W^5,W^6$
construction. The readable companion prints the complete nonlinear
rows and the seven-term decoupling relation so their algebraic
role can be checked separately. The finite-reconstruction theorem
then recovers the remaining generator brackets by an acyclic
sequence of Wick and iterate reductions. It is this combination,
not agreement of $(c,\lambda)$ alone, that identifies the PF13
factor.

\subsection{Reading old scope flags without reinstating obsolete gaps}

\begin{longtable}{@{}p{0.32\linewidth}p{0.29\linewidth}p{0.29\linewidth}@{}}
\toprule
Archived record&Original finite scope&Current use in this draft\\
\midrule\endhead
\artifact{hi_cft_20260923_proof_potts.json}&Potts stress exact;
weight-three membership numerical.&The separate exact primary and
membership record supplies that missing finite input.\\
\artifact{hi_cft_20260923_proof_pf13_membership.json}&Membership and
positive stress exact; PF identification explicitly false/conditional
in that one file.&The complete OPE, decoupling and reconstruction
chain supplies the identification separately. The old flag is
not edited.\\
\artifact{hi_cft_20260923_proof_core_coset.json}&Numerical dimensions
in selected $S_2,S_3$ only.&Retained as exploratory observations;
not promoted to a full coset character.\\
\artifact{hi_cft_20260923_next_structure_order13_core.md}&Actual
core established; no properness or $\Rep(U)$ identification.&Core
input to the top bound; properness uses the independent finite
strong-generation argument.\\
\artifact{hi_cft_20260923_next_structure_residual_low_character.md}&
Residual grades two and three conditional on properness and
$U_3=S_3$.&The present manuscript supplies those separate premises;
no claim about residual grades four and higher is added.\\
\artifact{hi_cft_20260923_next_structure_m6_character_obstruction.md}&
Two explicit ribbon models excluded; no general classification
claimed.&The later universal weighted projection supplies the
missing generality.\\
\artifact{hi_cft_20260923_next_structure_final_dependency_audit.md}&
Conditional identification coherent; properness and analytic bridges
not yet supplied at that date.&Useful audit of the algebraic chain;
current finite closure and analytic proofs are additional work.\\
\artifact{hi_cft_c8_revision_20260927_analytic.md}&Early entries
leave A1 open; final entry explicitly supersedes them.&Current
pairwise A1 proof is retained with its domain and lift obligations;
this appendix does not revert it to an assumption.\\
\bottomrule
\end{longtable}

A historical record can have a correct limited statement even
when later work proves more. Conversely, a later global argument
does not turn a numerical residual or a source hash into an
independent exact field verification. The source annotations above
preserve both directions of this distinction.

\subsection{Additional primary dependencies of the supplementary coset route}

The following source locations are retained from the archived notes;
the expansion does not claim a fresh external source audit.
They identify which theorem applications beyond the main proof a
reviewer should inspect when using the supplementary results.

\begin{longtable}{@{}p{0.27\linewidth}p{0.65\linewidth}@{}}
\toprule
Application&Primary source location recorded in the archive\\
\midrule\endhead
Lattice--parafermion reconstruction&Arakawa--Lam--Yamada
\cite{ALY}, Sections 2.4 and 6, equations (2.13), (2.16), (2.19)
and simple-current extension Theorem 2.4.\\
No proper extension of $R$&McRae,
\href{https://arxiv.org/abs/1810.00747}{compact-orbifold tensor theory},
Corollary 4.8 (and Theorem 4.14 for the optional regularity of a
hypothetical extension).\\
Index-twelve $F_4$ embedding&Penkov--Zuckerman,
\href{https://math.constructor.university/penkov/papers/gen_hc_r_final.pdf}{generalized Harish-Chandra modules},
Section 3 and Table B for the index list and uniqueness by index.\\
Finite centralizer $S_4$&Westaway,
\href{https://arxiv.org/abs/2312.01838}{Section 4.9}, and Xue,
\href{https://jolt.centre-mersenne.org/item/10.5802/jolt.950.pdf}{Section 5.1 and Table 13},
for the $F_4(a_3)$ diagram, nilpotent centralizer and component group.\\
Regular diagonal frame&Lin,
\href{https://arxiv.org/abs/2107.09400}{rationality of diagonal coset VOAs},
Theorem 5.7 at $k=8$, and Theorem 3.10 for the stated finite
extension $C_2$ step.\\
Ambient residual regularity&McRae,
\href{https://arxiv.org/abs/2108.01898}{Theorem 6.11}, after the
actual double-commutant and $C_2$ hypotheses have been checked.\\
Mirror branching&McRae,
\href{https://arxiv.org/abs/2107.06577}{mirror equivalence}, Theorem 1.1
and the twist discussion after Theorem 4.6.\\
Charge-six regularity&Creutzig--Kanade--Linshaw--Ridout,
\href{https://arxiv.org/abs/1611.00305}{Section 4.3}, including its
standing lattice hypothesis, Lemma 4.6, Theorem 4.12 and
Corollary 4.13.\\
Category splitting and actual extension transport&M\"uger modular
centralizer decomposition, and the Lagrangian/local-module and
extension-transitivity results used in \cite{DMNO,Mueger,HKL}.
The archived notes specify the diagonal algebra and nonzero dimension
rather than appealing only to a matching global dimension.\\
\bottomrule
\end{longtable}
